\section{AHSME 1962}
        \begin{enumerate}
        	\item (\textit{AHSME 1962}) The expression $\frac{1^{4y-1}}{5^{-1}+3^{-1}}$ is equal to:


 $\textbf{(A)}\ \frac{4y-1}{8} \qquad  \textbf{(B)}\ 8 \qquad  \textbf{(C)}\ \frac{15}{2} \qquad  \textbf{(D)}\ \frac{15}{8} \qquad  \textbf{(E)}\ \frac{1}{8}$

 \textbf{Answer: }D



\textbf{Solution}

We simplify the expression to yield:

 $\dfrac{1^{4y-1}}{5^{-1}+3^{-1}}=\dfrac{1}{5^{-1}+3^{-1}}=\dfrac{1}{\dfrac{1}{5}+\dfrac{1}{3}}=\dfrac{1}{\dfrac{8}{15}}=\dfrac{15}{8}$.

 Thus our answer is $\boxed{\textbf{(D)}\ \frac{15}{8}}$.

 


	\item (\textit{AHSME 1962}) The expression $\sqrt{\frac{4}{3}} - \sqrt{\frac{3}{4}}$ is equal to:


 $\textbf{(A)}\ \sqrt{3}/6 \qquad  \textbf{(B)}\ -\sqrt{3}/6 \qquad  \textbf{(C)}\ \sqrt{-3}/6 \qquad  \textbf{(D)}\ 5\sqrt{3}/6 \qquad  \textbf{(E)}\ 1$

 \textbf{Answer: }A



\textbf{Solution}

Simplifying $\sqrt{\dfrac{4}{3}}$ yields $\dfrac{2}{\sqrt{3}}=\dfrac{2\sqrt{3}}{3}$.

 
Simplifying $\sqrt{\dfrac{3}{4}}$ yields $\dfrac{\sqrt{3}}{2}$.

 
$\dfrac{2\sqrt{3}}{3}-\dfrac{\sqrt{3}}{2}=\dfrac{4\sqrt{3}}{6}-\dfrac{3\sqrt{3}}{6}=\dfrac{\sqrt{3}}{6}$.

 
Since we cannot simplify further, the correct answer is $\boxed{\textbf{(A)}\ \frac{\sqrt{3}}{6}}$



	\item (\textit{AHSME 1962}) The first three terms of an arithmetic progression are $x - 1, x + 1, 2x + 3$, in the order shown. The value of $x$ is:


 $\textbf{(A)}\ -2 \qquad  \textbf{(B)}\ 0 \qquad  \textbf{(C)}\ 2 \qquad  \textbf{(D)}\ 4 \qquad  \textbf{(E)}\ \text{undetermined}$

 \textbf{Answer: }B



\textbf{Solution}

Let $y$ represent the common difference between the terms. We have $(x+1)-y=(x-1)\implies y=2$.

 Substituting gives us $(2x+3)-2=(x+1)\implies 2x+1=x+1\implies x=0$.

 Therefore, our answer is $\boxed{\textbf{(B)}\ 0}$;

 


	\item (\textit{AHSME 1962}) If $8^x = 32$, then $x$ equals:


 $\textbf{(A)}\ 4 \qquad  \textbf{(B)}\ \frac{5}{3} \qquad  \textbf{(C)}\ \frac{3}{2} \qquad  \textbf{(D)}\ \frac{3}{5} \qquad  \textbf{(E)}\ \frac{1}{4}$

 \textbf{Answer: }B



\textbf{Solution}

Recognizing that $8=2^3$, we know that $2^{3x}=32$. Since $2^5=32$, we have $2^{3x}=2^5$. Therefore, $x=\dfrac{5}{3}$.

 So our answer is $\boxed{\textbf{(B)}\ \frac{5}{3}}$



	\item (\textit{AHSME 1962}) If the radius of a circle is increased by $1$ unit, the ratio of the new circumference to the new diameter is:


 $\textbf{(A)}\ \pi + 2 \qquad  \textbf{(B)}\ \frac{2\pi + 1}{2} \qquad  \textbf{(C)}\ \pi \qquad  \textbf{(D)}\ \frac{2\pi - 1}{2} \qquad  \textbf{(E)}\ \pi - 2$

 \textbf{Answer: }C



\textbf{Solution (Intuitive)}

The ratio of a circumference to a diameter always is the same so the answer is obviously C.

 



\textbf{Solution 2 (Full Proof)}

Let us say that the radius of a circle is $r$. When the radius is increased by $1$, the new radius is $r+1$ so the diameter is $2r+2$. We know that the circumference of a circle is $2\pi r$ so $2 \cdot \pi \cdot (r+1) = \pi \cdot (2r+2)$. Finally, the problem asked for the ratio of the new circumference to the new diameter is $\frac{\pi \cdot (2r+2)}{2r+2}=\boxed{\pi}$.

 


	\item (\textit{AHSME 1962}) A square and an equilateral triangle have equal perimeters. The area of the triangle is $9\sqrt{3}$ square inches. Expressed in inches the diagonal of the square is:


 $\textbf{(A)}\ 9/2 \qquad  \textbf{(B)}\ 2\sqrt{5} \qquad  \textbf{(C)}\ 4\sqrt{2} \qquad  \textbf{(D)}\ 9\sqrt{2}/2 \qquad  \textbf{(E)}\ \text{none of these}$

 \textbf{Answer: }D



\textbf{Solution}

To solve for the perimeter of the triangle we plug in the formula for the area of an equilateral triangle which is $\dfrac{x^2\sqrt{3}}{4}$. This has to be equal to $9 \sqrt{3}$, which means that $x=6$, or the side length of the triangle is $6$. Thus, the triangle (and the square) have a perimeter of $18$. It follows that each side of the square is $\dfrac{18}{4}=4.5$. If we draw the diagonal, we create a 45-45-90 triangle, whose hypotenuse (also the diagonal of the square) is $\boxed{\text{(D)} \dfrac{9\sqrt{2}}{2}}$.

 


	\item (\textit{AHSME 1962}) Let the bisectors of the exterior angles at $B$ and $C$ of triangle $ABC$ meet at $D$. Then the measure in degrees of angle $BDC$ is:


 $\textbf{(A)}\ \frac{1}{2}(90 - A) \qquad  \textbf{(B)}\ 90 - A \qquad  \textbf{(C)}\ \frac{1}{2}(180 - A) \qquad  \textbf{(D)}\ 180 - A \qquad  \textbf{(E)}\ 180 - 2A$

 \textbf{Answer: }C



\textbf{Solution}

Calculating for angles $\angle DBC$ and $\angle DCB$, we get

 $\angle DBC$ = $90 - \frac{B}{2}$ and$\angle DCB$ = $90 - \frac{C}{2}$.

 In triangle $BCD$, we have

 $\angle BDC$ = $180 - (90 - \frac{B}{2}) - (90 - \frac{C}{2})$ = $\frac{B+C}{2}$ = $\frac{1}{2}\cdot(180 - A)$, so the answer is $\fbox{C}$.

 


	\item (\textit{AHSME 1962}) Given the set of $n$ numbers, $n > 1$, of which one is $1 - (1/n)$, and all the others are $1$. The arithmetic mean of the $n$ numbers is:


 $\textbf{(A)}\ 1 \qquad  \textbf{(B)}\ n - \frac{1}{n} \qquad  \textbf{(C)}\ n - \frac{1}{n^2} \qquad  \textbf{(D)}\ 1 - \frac{1}{n^2} \qquad  \textbf{(E)}\ 1 - \frac{1}{n} - \frac{1}{n^2}$

 \textbf{Answer: }D



\textbf{Solution}

Just take $\frac{1(n-1)+(1-\frac{1}{n})}{n}$. You get $\frac{n-1+1-\frac{1}{n}}{n}$, which is just $\frac{n-\frac{1}{n}}{n}$, which is just $\boxed{D}$



	\item (\textit{AHSME 1962}) When $x^9 - x$ is factored as completely as possible into polynomials and monomials with real integral coefficients, the number of factors is:


 $\textbf{(A)}\ \text{more than } 5 \qquad  \textbf{(B)}\ 5 \qquad  \textbf{(C)}\ 4 \qquad  \textbf{(D)}\ 3 \qquad  \textbf{(E)}\ 2$

 \textbf{Answer: }B



\textbf{Solution}

Obviously, we can factor out an $x$ first to get $x(x^8-1)$.Next, we repeatedly factor differences of squares:
 \[x(x^4+1)(x^4-1)\]
 \[x(x^4+1)(x^2+1)(x^2-1)\]
 \[x(x^4+1)(x^2+1)(x+1)(x-1)\]None of these 5 factors can be factored further, so the answer is $\boxed{\textbf{(B) } 5}$.

 


	\item (\textit{AHSME 1962}) A man drives $150$ miles to the seashore in $3$ hours and $20$ minutes. He returns from the shore to the starting point in $4$ hours and $10$ minutes. Let $r$ be the average rate for the entire trip. Then the average rate for the trip going exceeds $r$, in miles per hour, by:


 $\textbf{(A)}\ 5 \qquad  \textbf{(B)}\ 4\frac{1}{2} \qquad  \textbf{(C)}\ 4 \qquad  \textbf{(D)}\ 2 \qquad  \textbf{(E)}\ 1$

 \textbf{Answer: }A



\textbf{Solution}

Since the equation for rate is $r=\frac{d}{t}$, you only need to find $d$ and $t$. The distance from the seashore to the starting point is $150$ miles, and since he makes a round trip, $d=300$. Also, you know the times it took him to go both directions, so when you add them up ($3$ hours and $20$ minutes $+$ $4$ hours and $10$ minutes), you get $7\frac{1}{2}$ hours. Since $r=\frac{d}{t}$, after plugging in the values of $d$ and $t$, you get $r=\frac{300}{7\frac{1}{2}}$, or $r=40$. But what the question asks is for the positive difference between the speed going to the seashore and the average speed. The equation is $r=\frac{d}{t}$, so when you plug in $d$ and $t$, you get $r=\frac{150}{3\frac{1}{3}}$. Simplifying, you get $r=45$. $45-40=5 \rightarrow \boxed{\text{A}}$



	\item (\textit{AHSME 1962}) The difference between the larger root and the smaller root of

 \[x^2 - px + (p^2 - 1)/4 = 0\]
is


 $\textbf{(A)}\ 0 \qquad  \textbf{(B)}\ 1 \qquad  \textbf{(C)}\ 2 \qquad  \textbf{(D)}\ p \qquad  \textbf{(E)}\ p + 1$

 \textbf{Answer: }B



\textbf{Solution}

Call the two roots $r$ and $s$, with $r \ge s$. By Vieta's formulas, $p=r+s$ and $(p^2-1)/4=rs.$(Multiplying both sides of the second equation by 4 gives $p^2-1=4rs$.)The value we need to find, then, is $r-s$.Since $p=r+s$, $p^2=r^2+2rs+s^2$.Subtracting $p^2-1=4rs$ from both sides gives $1=r^2-2rs+s^2$.Taking square roots, $r-s=1 \Rightarrow \boxed{\textbf{(B)}}$.(Another solution is to use the quadratic formula and see that theroots are $\frac{p\pm 1}2$, and their difference is 1.)

 


	\item (\textit{AHSME 1962}) When $\left(1 - \frac{1}{a}\right)^6$ is expanded, the sum of the last three coefficients is:


 $\textbf{(A)}\ 22 \qquad  \textbf{(B)}\ 11 \qquad  \textbf{(C)}\ 10 \qquad  \textbf{(D)}\ -10 \qquad  \textbf{(E)}\ -11$

 \textbf{Answer: }C



\textbf{Solution}

This is equivalent to $\frac{(a-1)^6}{a^6}.$Its expansion has 7 terms, whose coefficients are the same as those of $(a-1)^6$.By the Binomial Theorem, the sum of the last three coefficients is $\binom{6}{2}-\binom{6}{1}+\binom{6}{0}=15-6+1=\boxed{10 \textbf{ (C)}}$.

 


	\item (\textit{AHSME 1962}) $R$ varies directly as $S$ and inversely as $T$. When $R = \frac{4}{3}$ and $T = \frac{9}{14}$, $S = \frac{3}{7}$. Find $S$ when $R = \sqrt{48}$ and $T = \sqrt{75}$.


 $\textbf{(A)}\ 28 \qquad  \textbf{(B)}\ 30 \qquad  \textbf{(C)}\ 40 \qquad  \textbf{(D)}\ 42 \qquad  \textbf{(E)}\ 60$

 \textbf{Answer: }B



\textbf{Solution}


 \[R=c\cdot\frac{S}T\]

 for some constant $c$. 

 You know that

 
 \[\frac43=c\cdot\frac{3/7}{9/14}=c\cdot\frac37\cdot\frac{14}9=c\cdot\frac23\,,\]

 so

 
 \[c=\frac{4/3}{2/3}=2\,.\]

 When $R=\sqrt{48}$ and $T=\sqrt{75}$ we have

 
 \[\sqrt{48}=\frac{2S}{\sqrt{75}}\,,\]

 so

 
 \[S=\frac12\sqrt{48\cdot75}=30\,.\] $\boxed{B}$
\begin{verbatim}
-- zixuan 12
\end{verbatim}




	\item (\textit{AHSME 1962}) Let $s$ be the limiting sum of the geometric series $4 - \frac{8}{3} + \frac{16}{9} - \cdots$, as the number of terms increases without bound. Then $s$ equals:


 $\textbf{(A)}\ \text{a number between } 0 \text{ and } 1 \qquad  \textbf{(B)}\ 2.4 \qquad  \textbf{(C)}\ 2.5 \qquad  \textbf{(D)}\ 3.6 \qquad  \textbf{(E)}\ 12$

 \textbf{Answer: }B



\textbf{Solution}

The infinite sum of a geometric series with first term $a$ and common ratio $r$ ($-1<r<1$) is $\frac{a}{1-r}$.Now, in this geometric series, $a=4$, and $r=-\frac23$. Plugging these into the formula, we get $\frac4{1-(-\frac23)}$, which simplifies to $\frac{12}5$, or $\boxed{2.4\textbf{ (B)}}$.

 


	\item (\textit{AHSME 1962}) Given triangle $ABC$ with base $AB$ fixed in length and position. As the vertex $C$ moves on a straight line, the intersection point of the three medians moves on:


 $\textbf{(A)}\ \text{a circle} \qquad  \textbf{(B)}\ \text{a parabola} \qquad  \textbf{(C)}\ \text{an ellipse} \qquad  \textbf{(D)}\ \text{a straight line} \qquad  \textbf{(E)}\ \text{a curve not listed here}$

 \textbf{Answer: }D



\textbf{Solution}

Let $CM$ be the median through vertex $C$, and let $G$ be the point of intersection of the triangle's medians. 

 Let $CH$ be the altitude of the triangle through vertex $C$ and $GP$ be the distance from $G$ to $AB$, with the point $P$ laying on $AB$. 

 Using Thales' intercept theorem, we derive the proportion:

 
$\frac{GP}{CH} = \frac{GM}{CM}$

 
The fraction $\frac{GM}{CM}$ in any triangle is equal to $\frac{1}{3}$ . Therefore $GP = \frac{CH}{3}$ . 

 
Since the problem states that the vertex $C$ is moving on a straight line, the length of $CH$ is a constant value. That means that the length of $GP$ is also a constant. Therefore the point $G$ is moving on a straight line.

 Answer: D

 


	\item (\textit{AHSME 1962}) Given rectangle $R_1$ with one side $2$ inches and area $12$ square inches. Rectangle $R_2$ with diagonal $15$ inches is similar to $R_1$. Expressed in square inches the area of $R_2$ is:


 $\textbf{(A)}\ 9/2 \qquad  \textbf{(B)}\ 36 \qquad  \textbf{(C)}\ 135/2 \qquad  \textbf{(D)}\ 9\sqrt{10} \qquad  \textbf{(E)}\ 27\sqrt{10}/4$

 \textbf{Answer: }C



\textbf{Solution}

Clearly, the other side of $R_1$ has length $\frac{12}2=6$.Now call the sides of $R_2$ a and b, with $b>a$. We know that $\frac{b}a=\frac31=3$, because the two rectangles are similar. We also know that $a^2+b^2=15^2=225$. But $b=3a$, so substituting gives 
 \[a^2+(3a)^2=225\]
 \[10a^2=225\]
 \[a^2=\frac{45}2\]
 \[a=\frac{3\sqrt{10}}2\]
 \[b=\frac{9\sqrt{10}}2\]
 \[[R_2]=ab=\boxed{\frac{135}2\textbf{ (C)}}\]



	\item (\textit{AHSME 1962}) If $a = \log_8 225$ and $b = \log_2 15$, then $a$, in terms of $b$, is:


 $\textbf{(A)}\ b/2 \qquad  \textbf{(B)}\ 2b/3 \qquad  \textbf{(C)}\ b \qquad  \textbf{(D)}\ 3b/2 \qquad  \textbf{(E)}\ 2b$

 \textbf{Answer: }B



\textbf{Solution}

Using the change-of-base rule: $a = \frac{\log 225}{\log 8}$ and $b = \frac{\log 15}{\log 2}$.
 \[\frac{a}{b} = \frac{\log 225 \log 2}{\log 8 \log 15}\]
 \[a = b \cdot \frac{\log 225}{\log 15} \cdot \frac{\log 2}{\log 8}\]
 \[a = b \log_{15} 225 \log_8 2\]
 \[a = \boxed{\frac{2b}3 \textbf{ (B)}}\]



	\item (\textit{AHSME 1962}) A regular dodecagon ($12$ sides) is inscribed in a circle with radius $r$ inches. The area of the dodecagon, in square inches, is:


 $\textbf{(A)}\ 3r^2 \qquad  \textbf{(B)}\ 2r^2 \qquad  \textbf{(C)}\ 3r^2\sqrt{3}/4 \qquad  \textbf{(D)}\ r^2\sqrt{3} \qquad  \textbf{(E)}\ 3r^2\sqrt{3}$

 \textbf{Answer: }A



\textbf{Solution}

The formula for the area of a regular dodecagon is $3r^2$. The answer is $\boxed{\textbf{(A)}}$.(If you don't know this formula, it's pretty easy to figure out that the area of a square inscribed in a circle is $2r^2$, and all the choices except $3r^2$ are less than $2r^2$. Remember, the more sides a regular polygon has, the closer its area gets to $\pi r^2$.)

 


	\item (\textit{AHSME 1962}) If the parabola $y = ax^2 + bx + c$ passes through the points $(-1, 12)$, $(0, 5)$, and $(2, -3)$, the value of $a + b + c$ is:


 $\textbf{(A)}\ -4 \qquad  \textbf{(B)}\ -2 \qquad  \textbf{(C)}\ 0 \qquad  \textbf{(D)}\ 1 \qquad  \textbf{(E)}\ 2$

 \textbf{Answer: }C



\textbf{Solution}

Substituting in the $(x, y)$ pairs gives the following system of equations:
 \[a-b+c=12\]
 \[c=5\]
 \[4a+2b+c=-3\]We know $c=5$, so plugging this in reduces the system to two variables:
 \[a-b=7\]
 \[4a+2b=-8\]Dividing the second equation by 2 gives $2a+b=-4$, which can be added to the first equation to get $3a=3$, or $a=1$. So the solution set is $(1, -6, 5)$, and the sum is $\boxed{0\textbf{ (C)}}$.

 


	\item (\textit{AHSME 1962}) The angles of a pentagon are in arithmetic progression. One of the angles, in degrees, must be:


 $\textbf{(A)}\ 108 \qquad  \textbf{(B)}\ 90 \qquad  \textbf{(C)}\ 72 \qquad  \textbf{(D)}\ 54 \qquad  \textbf{(E)}\ 36$

 \textbf{Answer: }A



\textbf{Solution}

If the angles are in an arithmetic progression, they can be expressed as$a$, $a+n$, $a+2n$, $a+3n$, and $a+4n$ for some real numbers $a$ and $n$.Now we know that the sum of the degree measures of the angles of a pentagon is $180(5-2)=540$.Adding our expressions for the five angles together, we get $5a+10n=540$.We now divide by 5 to get $a+2n=108$. It so happens that $a+2n$ is one of the angles we defined earlier, so that angle must have a measure of $\boxed{108\textbf{ (A)}}$.(In fact, for any arithmetic progression with an odd number of terms,the middle term is equal to the average of all the terms.)

 



\textbf{Solution 2}

If we write the five terms as $a$, $a - n$, $a - 2n$, $a + n$ and $a + 2n$, we can see that adding them up, we get $5a = 540$ through this, we can see that $a = 108$, $\fbox{\textbf{(A)}}$



	\item (\textit{AHSME 1962}) It is given that one root of $2x^2 + rx + s = 0$, with $r$ and $s$ real numbers, is $3 + 2i$ ($i = \sqrt{-1}$). The value of $s$ is:


 $\textbf{(A)}\ \text{undetermined} \qquad  \textbf{(B)}\ 5 \qquad  \textbf{(C)}\ 6 \qquad  \textbf{(D)}\ -13 \qquad  \textbf{(E)}\ 26$

 \textbf{Answer: }E



\textbf{Solution}

If a quadratic with real coefficients has two non-real roots, the two roots must be complex conjugates of one another.This means the other root of the given quadratic is $\overline{3+2i}=3-2i$. Now Vieta's formulas say that $s/2$ is equal to the product of the two roots, so$s = 2(3+2i)(3-2i) = \boxed{26 \textbf{ (E)}}$.

 


	\item (\textit{AHSME 1962}) The number $121_b$, written in the integral base $b$, is the square of an integer for


 $\textbf{(A)}\ b = 10,\text{ only} \qquad  \textbf{(B)}\ b = 10 \text{ and } b = 5,\text{ only} \qquad  \textbf{(C)}\ 2 \leq b \leq 10 \qquad  \textbf{(D)}\ b > 2 \qquad  \textbf{(E)}\ \text{no value of } b$

 \textbf{Answer: }D



\textbf{Solution}

$121_b$ can be represented in base 10 as $b^2+2b+1$, which factors as $(b+1)^2$.Note that $b>2$ because 2 is a digit in the base-b representation, but for any $b>2$, $121_b$ is the square of $b+1$. $\boxed{\textbf{(D)}}$



	\item (\textit{AHSME 1962}) In triangle $ABC$, $CD$ is the altitude to $AB$, and $AE$ is the altitude to $BC$. If the lengths of $AB$, $CD$, and $AE$ are known, the length of $DB$ is:


 $\textbf{(A)}\ \text{not determined by the information given} \qquad  \textbf{(B)}\ \text{determined only if } A \text{ is an acute angle} \qquad  \textbf{(C)}\ \text{determined only if } B \text{ is an acute angle} \qquad  \textbf{(D)}\ \text{determined only if } ABC \text{ is an acute triangle} \qquad  \textbf{(E)}\ \text{none of these is correct}$

 \textbf{Answer: }E



\textbf{Solution}

We can actually determine the length of $DB$ no matter what type of angles $A$ and $B$ are. This can be easily proved through considering all possible cases, though for the purposes of this solution, we'll show that we can determine $DB$ if $A$ is an obtuse angle. 

 Let's see what happens when $A$ is an obtuse angle. $\triangle AEB\sim \triangle CDB$ by AA, so $\frac{AE}{CD}=\frac{AB}{DB}$. Hence $DB=\frac{AE}{CD\times AB}$. Since we've determined the length of $DB$ even though we have an obtuse angle, $DB$ is not $\textbf{only}$ determined by what type of angle $A$ may be. Hence our answer is $\fbox{E}$.

 


	\item (\textit{AHSME 1962}) Three machines $P$, $Q$, and $R$, working together, can do a job in $x$ hours. When working alone $P$ needs an additional $6$ hours to do the job; $Q$, one additional hour; and $R$, $x$ additional hours. The value of $x$ is:


 $\textbf{(A)}\ \frac{2}{3} \qquad  \textbf{(B)}\ \frac{11}{18} \qquad  \textbf{(C)}\ \frac{3}{2} \qquad  \textbf{(D)}\ 2 \qquad  \textbf{(E)}\ 3$

 \textbf{Answer: }A



\textbf{Solution}

Machine P takes $x+6$ hours, machine Q takes $x+1$ hours, and machine R takes $2x$ hours.We also know that all three working together take $x$ hours. Now the time it takes for all three machines to complete the job is the harmonic mean of the times the three machines take individually; that is,
 \[x=\frac1{\frac1{x+6}+\frac1{x+1}+\frac1{2x}}\]
 \[\frac{x}{x+6}+\frac{x}{x+1}+\frac{x}{2x}=1\]
 \[\frac{x}{x+6}+\frac{x}{x+1}=\frac12\]
 \[2x(x+1)+2x(x+6)=(x+1)(x+6)\]
 \[2x^2+2x+2x^2+12x=x^2+7x+6\]
 \[3x^2+7x-6=0\]
 \[(3x-2)(x+3)=0\]
 \[x\in\{\frac23, -3\}\]Obviously, the number of hours is positive, so the answer is$\boxed{\frac23 \textbf{ (A)}}$.

 


	\item (\textit{AHSME 1962}) Given square $ABCD$ with side $8$ feet. A circle is drawn through vertices $A$ and $D$ and tangent to side $BC$. The radius of the circle, in feet, is:


 $\textbf{(A)}\ 4 \qquad  \textbf{(B)}\ 4\sqrt{2} \qquad  \textbf{(C)}\ 5 \qquad  \textbf{(D)}\ 5\sqrt{2} \qquad  \textbf{(E)}\ 6$

 \textbf{Answer: }C



\textbf{Solution}

Let $O$ be the center of the circle and $E$ be the point of tangency of the circle and $BC$ and let $F$ be the point of intersection of lines $OE$ and $AD$ Because of the symmetry, $BE=EC=AF=FD=4$ feet. Let the length of $OF$ be $x$. The length of $OE$ is $EF-OF=-x+8$. By Pythagorean Theorem, $OA=OD=\sqrt{x^2+4^2}=\sqrt{x^2+16}$. Because $OA$, $OD$, and $OE$ are radii of the same circle, $-x+8=OE=OA=AD=\sqrt{x^2+16}$. So, $\sqrt{x^2+16}=-x+8$. Squaring both sides, we obtain $x^2+16=x^2-16x+64$. Subtracting $x^2+16$ from both sides and adding $16x$, our equation becomes $16x=80$, so our answer is $x=\boxed{C)5}$.

 


	\item (\textit{AHSME 1962}) For any real value of $x$ the maximum value of $8x - 3x^2$ is:


 $\textbf{(A)}\ 0 \qquad  \textbf{(B)}\ \frac{8}{3} \qquad  \textbf{(C)}\ 4 \qquad  \textbf{(D)}\ 5 \qquad  \textbf{(E)}\ \frac{16}{3}$

 \textbf{Answer: }E



\textbf{Solution}

Let $f(x) = 8x-3x^2$ Since $f(x)$ is a quadratic and the quadratic term is negative, the maximum will be $f\bigg(- \dfrac{b}{2a}\bigg)$ when written in the form $ax^2+bx+c$. We see that $a=-3$, and so $- \dfrac{b}{2a} = -\bigg( \dfrac{8}{-6}\bigg) = \dfrac{4}{3}$. Plugging in this value yields $f(\dfrac{4}{3}) = \dfrac{32}{3}-3 \cdot \dfrac{16}{9} = \dfrac{32}{3} - \dfrac{16}{3} = \boxed{\text{E}\ \dfrac{16}{3}}$



	\item (\textit{AHSME 1962}) Let $a \textcircled{L} b$ represent the operation on two numbers, $a$ and $b$, which selects the larger of the two numbers, with $a \textcircled{L} a = a$. Let $a \textcircled{S} b$ represent the operation which selects the smaller of the two numbers, with $a \textcircled{S} a = a$. Which of the following three rules is (are) correct?


 (1) $a \textcircled{L} b = b \textcircled{L} a$, (2) $a \textcircled{L} (b \textcircled{L} c) = (a \textcircled{L} b) \textcircled{L} c$,
(3) $a \textcircled{S} (b \textcircled{L} c) = (a \textcircled{S} b) \textcircled{L} (a \textcircled{S} c)$.


 $\textbf{(A)}\ (1)\text{ only} \qquad  \textbf{(B)}\ (2)\text{ only} \qquad  \textbf{(C)}\ (1)\text{ and }(2)\text{ only} \qquad  \textbf{(D)}\ (1)\text{ and }(3)\text{ only} \qquad  \textbf{(E)}\ \text{all three}$

 \textbf{Answer: }E



\textbf{Solution}

The first rule must be correct as both sides of the equation pick the larger out of a and b.

 The second rule must also be correct as both sides would end up picking the largest out of a, b, and c.

 WLOG, lets assume b < c.

 The third rule is a little more complex. To see if it works, let’s split the possibilities into three cases. a < b < c, b < a < c, and b < c < a. For the first case, the equation simplifies to a = a, which is correct. In the second case, the equation simplifies to a = a, which is correct. For the last case, the equation simplifies to c = c, which is also correct. Since all three cases work, this rule also works.

 Thus, our answer is $\boxed{\textbf{E}}$

 (Solution by someonenumber011)

 


	\item (\textit{AHSME 1962}) The set of $x$-values satisfying the equation $x^{\log_{10} x} = x^3/100$ consists of:


 $\textbf{(A)}\ \frac{1}{10},\text{ only} \qquad  \textbf{(B)}\ 10,\text{ only} \qquad  \textbf{(C)}\ 100,\text{ only} \qquad  \textbf{(D)}\ 10 \text{ or } 100,\text{ only} \qquad  \textbf{(E)}\ \text{more than two real numbers.}$

 \textbf{Answer: }D



\textbf{Solution}

Taking the base-$x$ logarithm of both sides gives $\log_{10}x=\log_x\frac{x^3}{100}$.This simplifies to 
 \[\log_{10}x=\log_x{x^3} - \log_x{100}\]
 \[\log_{10}x+\log_x{100}=3\]
 \[\log_{10}x+2 \log_x{10}=3\]At this point, we substitute $u=\log_{10}x$.
 \[u+\frac2u=3\]
 \[u^2+2=3u\]
 \[u^2-3u+2=0\]
 \[(u-2)(u-1)=0\]
 \[u\in\{1, 2\}\]
 \[x\in\{10, 100\}\]The answer is $\boxed{\textbf{(D)}}$.

 


	\item (\textit{AHSME 1962}) Which of the following sets of $x$-values satisfy the inequality

 \[2x^2 + x < 6?\]
$\textbf{(A)}\ -2 < x < \frac{3}{2} \qquad  \textbf{(B)}\ x > \frac{3}{2} \text{ or } x < -2 \qquad  \textbf{(C)}\ x < \frac{3}{2} \qquad  \textbf{(D)}\ \frac{3}{2} < x < 2 \qquad  \textbf{(E)}\ x < -2$

 \textbf{Answer: }A



\textbf{Solution}

First, subtract 6 from both sides of the inequality, $2x^2 + x - 6 < 0$.This is a parabola that opens upward when graphed, it has a positive leading coefficient. So any negative x values must be between its x-axis intersections, namely $x = -2, 1.5$. The answer is A.

 


	\item (\textit{AHSME 1962}) Form I


 Consider the statements: (1) $p \wedge q$ (2) $p \wedge \sim q$ (3) $\sim p \wedge q$ (4) $\sim p \wedge \sim q$, where $p$ and $q$ are statements.


 How many of these imply the truth of $\sim(p \wedge q)$?


 Form II


 Consider the statements: (1) $p$ and $q$ are both true (2) $p$ is true and $q$ is false (3) $p$ is false and $q$ is true (4) $p$ is false and $q$ is false.


 How many of these imply the negation of the statement "$p$ and $q$ are both true"?


 $\textbf{(A)}\ 0 \qquad  \textbf{(B)}\ 1 \qquad  \textbf{(C)}\ 2 \qquad  \textbf{(D)}\ 3 \qquad  \textbf{(E)}\ 4$

 \textbf{Answer: }D



\textbf{Solution}

By De Morgan's Law, the negation of $p$ and $q$ are both true is that at least one of them is false, with the exception of the statement 1, i.e.;

 $\text{p and q are both true}.$

 The other three statements state that at least one statement is false. So, 2, 3, and 4 work, yielding an answer of 3, or $\textbf{(D)}.$



	\item (\textit{AHSME 1962}) The ratio of the interior angles of two regular polygons is $3:2$. How many such pairs are there?


 $\textbf{(A)}\ 1 \qquad  \textbf{(B)}\ 2 \qquad  \textbf{(C)}\ 3 \qquad  \textbf{(D)}\ 4 \qquad  \textbf{(E)}\ \text{infinitely many}$

 \textbf{Answer: }C



\textbf{Solution}

The formula for the measure of the interior angle of a regular polygon with $n$-sides is $180 - \frac{360}{n}$. Letting our two polygons have side length $r$ and $k$, we have that the ratio of the interior angles is $\frac{180 - \frac{360}{r}}{180 - \frac{360}{k}} = \frac{(r-2) \cdot k}{(k-2) \cdot r} = \frac{3}{2}$.  Cross multiplying both sides, we have $2rk-4k = 3kr-6r \Rightarrow -rk-4k+6r = 0$. Using Simon's Favorite Factoring Trick, we have $-k(r+4)+6r+24 = 24 \Rightarrow (r+4)(6-k)=24$. Because $k$ and $r$ are both more than $2$, we know that $6-k < r+4$. Now, we just set these factors equal to the factors of 24. We can set $6-k$ to $1$, $2$, $3$, or $4$ and $r+4$ to $24$, $12$, $8$, or $6$ respectively to get the following pairs for $(k, r)$: $(5, 20)$, $(4, 8)$, $(3, 4)$, and $(2, 2)$. However, we have to take out the solution with $k = 2$, because $k$ and $r$ are both more than $2$, leaving us with $\boxed{\textbf{(C)}\ 3}$ as the correct answer.

 


	\item (\textit{AHSME 1962}) If $x_{k+1} = x_k + \frac{1}{2}$ for $k = 1, 2, \cdots, n - 1$ and $x_1 = 1$, find

 \[x_1 + x_2 + \cdots + x_n.\]
$\textbf{(A)}\ \frac{n+1}{2} \qquad  \textbf{(B)}\ \frac{n+3}{2} \qquad  \textbf{(C)}\ \frac{n^2-1}{2} \qquad  \textbf{(D)}\ \frac{n^2+n}{4} \qquad  \textbf{(E)}\ \frac{n^2+3n}{4}$

 \textbf{Answer: }E



\textbf{Solution}

The sequence $x_1, x_2, \dots, x_n$ is an arithmetic sequence since every term is $\frac12$ more than the previous term. Letting $a=1$ and $r=\frac12$, we can rewrite the sequence as $a, a+r, \dots, a+(n-1)r$.Recall that the sum of the first $n$ terms of an arithmetic sequence is $na+\binom{n}2r$.Substituting our values for $a$ and $r$, we get $n+\frac{\binom{n}2}{2}$. Simplifying gives$\boxed{\frac{n^2+3n}{4}\textbf{ (E)}}$



	\item (\textit{AHSME 1962}) The set of $x$-values satisfying the inequality $2 \leq |x - 1| \leq 5$ is:


 $\textbf{(A)}\ -4 \leq x \leq -1 \text{ or } 3 \leq x \leq 6 \qquad  \textbf{(B)}\ 3 \leq x \leq 6 \text{ or } -6 \leq x \leq -3 \qquad  \textbf{(C)}\ x \leq -1 \text{ or } x \geq 3 \qquad  \textbf{(D)}\ -1 \leq x \leq 3 \qquad  \textbf{(E)}\ -4 \leq x \leq 6$

 \textbf{Answer: }A



\textbf{Solution}

This inequality can be split into two inequalities:$2\le x-1\le5$ or $2\le1-x\le5$. Solving for x gives$3\le x\le6$ or $-4\le x\le-1$. $\boxed{\textbf{(A)}}$



	\item (\textit{AHSME 1962}) For what real values of $K$ does $x = K^2(x - 1)(x - 2)$ have real roots?


 $\textbf{(A)}\ \text{none} \qquad  \textbf{(B)}\ -2 < K < 1 \qquad  \textbf{(C)}\ -2\sqrt{2} < K < 2\sqrt{2} \qquad  \textbf{(D)}\ K > 1 \text{ or } K < -2 \qquad  \textbf{(E)}\ \text{all}$

 \textbf{Answer: }E



\textbf{Solution}

Factoring, we get
 \[x = K^2x^2 - 3K^2x + 2K^2\]
 \[K^2x^2 - (3K^2+1)x + 2K^2 = 0\]At this point, we could use the quadratic formula, but we're only interested in whether the roots are real, which occurs if and only if the discriminant ($b^2-4ac$) is non-negative.
 \[(-3K^2-1)^2-8K^4\ge0\]
 \[9K^4+6K^2+1-8K^4\ge0\]
 \[K^4+6K^2+1\ge0\]This function of $K$ is symmetric with respect to the y-axis. Plugging in $0$ for $K$ gives us $1$, and it is clear that the function is strictly increasing over $(0,\infty)$, and so it is positive for all real values of $K$. $\boxed{\textbf{(E)}}$



	\item (\textit{AHSME 1962}) A man on his way to dinner shortly after $6{:}00$ p.m. observes that the hands of his watch form an angle of $110^\circ$. Returning before $7{:}00$ p.m. he notices that again the hands of his watch form an angle of $110^\circ$. The number of minutes that he has been away is:


 $\textbf{(A)}\ 36\frac{2}{3} \qquad  \textbf{(B)}\ 40 \qquad  \textbf{(C)}\ 42 \qquad  \textbf{(D)}\ 42.4 \qquad  \textbf{(E)}\ 45$

 \textbf{Answer: }B



\textbf{Solution}

Let $n$ be the number of minutes after 6:00. Let $h(n)=180+\frac{n}2$ be the angle, in degrees, of the hour hand (with $0^{\circ}$ at the top and increasing in the clockwise direction); similarly, let $m(n)=6n$ be the angle of the minute hand. We want $|h(n)-m(n)|=110$. This is equivalent to 
 \[180-\frac{11n}2=\pm110\]
 \[-\frac{11n}2\in\{-70,-290\}\]
 \[\frac{11n}2\in\{70,290\}\]
 \[11n\in\{140,580\}\]
 \[n\in\{\frac{140}{11},\frac{580}{11}\}\]The difference between the two values of $n$ is $\frac{440}{11}=\boxed{40\textbf{ (B)}}$.

 


	\item (\textit{AHSME 1962}) If both $x$ and $y$ are integers, how many solutions are there of the equation $(x - 8)(x - 10) = 2^y$?


 $\textbf{(A)}\ 0 \qquad  \textbf{(B)}\ 1 \qquad  \textbf{(C)}\ 2 \qquad  \textbf{(D)}\ 3 \qquad  \textbf{(E)}\ \text{more than } 3$

 \textbf{Answer: }C



\textbf{Solution}

The equality implies $x-8$ and $x-10$ are both powers of two; since they differ by two, it must be the case that $(x-8,x-10) = (4,2)$ or $(x-8,x-10) = (-2,-4)$.  (Note that $(1,-1)$ is not allowed because then the product is negative.) These yield $(x,y) = (12,3)$ or $(x,y) = (6,3)$, for a total of $\boxed{2\textbf{ (C)}}$ solutions.

 


	\item (\textit{AHSME 1962}) $ABCD$ is a square with side of unit length. Points $E$ and $F$ are taken respectively on sides $AB$ and $AD$ so that $AE = AF$ and the quadrilateral $CDFE$ has maximum area. In square units this maximum area is:


 $\textbf{(A)}\ \frac{1}{2} \qquad  \textbf{(B)}\ \frac{9}{16} \qquad  \textbf{(C)}\ \frac{19}{32} \qquad  \textbf{(D)}\ \frac{5}{8} \qquad  \textbf{(E)}\ \frac{2}{3}$

 \textbf{Answer: }D



\textbf{Solution}

Let $AE=AF=x$$[CDFE]=[ABCD]-[AEF]-[EBC]=1-\frac{x^2}{2}-\frac{1-x}{2}$Or$[CDFE]=\frac{\frac{5}{4}-(x-\frac{1}{2})^2}{2}\le \frac{5}{8}$As $(x-\frac{1}{2})^2\ge 0$So $[CDFE]\le \frac{5}{8}$Equality occurs when $AE=AF=x=\frac{1}{2}$So the maximum value is $\frac{5}{8}$




\textbf{Solution 2}

Let us first draw a unit square.
                    \begin{figure}[H]
                        \centering
                        \includegraphics[width=0.5\textwidth]{Solution/AHSME_I_1962_37_Solution_Figure_1.png}
                    \end{figure}
                    We will now pick arbitrary points $E$ and $F$ on $\overline{AB}$ and $\overline{AD}$ respectively. We shall say that $AE = AF = x$
                    \begin{figure}[H]
                        \centering
                        \includegraphics[width=0.5\textwidth]{Solution/AHSME_I_1962_37_Solution_Figure_2.png}
                    \end{figure}
                    Thus, our problem has been simplified to maximizing the area of the blue quadrilateral.If we drop an altitude from $E$ to $\overline{DC}$, and call the foot of the altitude $G$, we can find the area of $EFDC$ by noting that $[EFDC] = [EGDF] + [EGC]$.
                    \begin{figure}[H]
                        \centering
                        \includegraphics[width=0.5\textwidth]{Solution/AHSME_I_1962_37_Solution_Figure_3.png}
                    \end{figure}
                    We can now finish the problem.

 Since $EG = 1$ and $FD = 1-x$, we have:
 \[[EFDC] = [EGDF] + [EGC]  = DG\cdot\frac{EG + FD}{2} + \frac{(CG)(EG)}{2}  = \frac{x(2-x)}{2} + \frac{1-x}{2}  = \frac{1+x-x^2}{2}\]

 To maximize this, we compete the square in the numerator to have:
 \[[EFDC] = \frac{-(x-\frac{1}{2})^2 + \frac{5}{4}}{2}  = \frac{5}{8} - \frac{(x-\frac{1}{2})^2}{2}\]

 Finally, we see that $\frac{(x-\frac{1}{2})^2}{2}\geq0$, as:
 \[\left(x-\frac{1}{2}\right)^2\geq0\]So,
 \[\frac{(x-\frac{1}{2})^2}{2}\geq0,\]

 where the first inequality was from the Trivial Inequality, and the second came from dividing both sides by $2$, which does not change the inequality sign.Thus, the maximum area is $\boxed{\frac{5}{8}}$ or $\boxed{\textbf{(D)}},$ when $\frac{(x-\frac{1}{2})^2}{2} = 0.$




\textbf{Solution 3 (Calculus)}

Let $AE = AF = x$. 

 The area of the quadrilateral $CDFE$ can be expressed as $A = 1 - \frac{x^2}{2} - \frac{(1-x)}{2}$.

 By taking the derivative of $A$, we get $A' = -x + \frac{1}{2}$.

 We make $A' = 0$ and get the critical point $x = \frac{1}{2}$.

 Substituting $x = \frac{1}{2}$ , the maximum area is $\boxed{\frac{5}{8}}$ or $\boxed{\textbf{(D)}}$.

 



\textbf{Solution 4}

Plot the points of the quadrilateral (with the bottom left corner of the square being at the origin) to plug into shoelace formula:
 \[(0, 0), (0, 1-x), (x, 1), (1, 0)\]So, 
 \[\frac{-x^2+x+1}{2}\]Trying $1/2$ first, you get $5/8$. Any larger or smaller value gives an answer less than $5/8$. Therefore, the answer is $\boxed{D}$.

 


	\item (\textit{AHSME 1962}) The population of Nosuch Junction at one time was a perfect square. Later, with an increase of $100$, the population was one more than a perfect square. Now, with an additional increase of $100$, the population is again a perfect square. The original population is a multiple of:


 $\textbf{(A)}\ 3 \qquad  \textbf{(B)}\ 7 \qquad  \textbf{(C)}\ 9 \qquad  \textbf{(D)}\ 11 \qquad  \textbf{(E)}\ 17$

 \textbf{Answer: }B



\textbf{Solution 1}

Let $a^2$ $=$ original population count, $b^2+1$ $=$ the second population count, and $c^2$ $=$ the third population count We first see that $a^2 + 100 = b^2 + 1$ or $99$ $=$ $b^2-a^2$.We then factor the right side getting $99$ $=$ $(b-a)(b+a)$. Since we can only have an nonnegative integral population, clearly $b+a$ $>$ $b-a$ and both factor $99$.We factor $99$ into $3^2 \cdot 11$ $=$ $(b-a)(b+a)$There are a few cases to look at:$1)$ $b+a$ $=$ $11$ and $b-a$ $=$ $9$.Adding the two equations we get $2b$ $=$ $20$ or $b$ $=$ $10$, which means $a$ $=$ $1$.But looking at the restriction that the second population + $100$ $=$ third population...$10^2$ $+$ $1$ $+$ $100$ $=$ $201$ $\neq$ a perfect square.

 $2)$ $b+a$ $=$ $33$ and $b-a$ $=$ $3$.Adding the two equations we get $2b$ $=$ $36$ or $b$ $=$ $18$, which means $a$ $=$ $15$.Looking at the same restriction, we get $18^2$ + $1$ + $100$ $=$ $324$ + $101$ $=$ $425$, which is NOT a perfect square.

 Finally, $b+a$ $=$ $99$ and $b-a$ $=$ $1$.$2b$ $=$ $100$ or $b$ $=$ $50$, which means $a$ $=$ $49$.Looking at the same restriction, we get $50^2$ + $1$ + $100$ $=$ $2500$ + $101$ $=$ $2601$ $=$ $51^2$. Thus we find that the original population is $a^2$ $=$ $49^2$ $=$ $7^4$. Or $a^2$ is a multiple of $\boxed{ (B) 7}$




\textbf{Solution 2}

Let $P$ $=$ original population. Translating the word problem into a system of equations, we got:
 \begin{align} P &= x^2 \\ P + 100 &= y^2 + 1 \\ P + 200 &= z^2 \end{align}for some positive integers $x$, $y$ and $z$.Now, by subtracting $(2)$ from $(3)$ (i.e. $(3) - (2)$), we got:
 \begin{align*} 100 &= z^2 - y^2 - 1 \\ 101 &= z^2 - y^2 \\ 101 &= (z - y)(z + y) \end{align*}
Since y and z are both positive integers and 101 is a prime, by factoring, the only working solution for us is $y = 50$ and $z = 51$.Plugging that back to $(2)$ and simplify, we got $P = 2401 = (49)^2 = x^2$, a multiple of $7$.Therefore, the answer is $\boxed{(B) 7}$. -nullptr07

 


	\item (\textit{AHSME 1962}) The medians $AN$ and $BP$ of a triangle with unequal sides are, respectively, $3$ inches and $6$ inches long. Its area is $3\sqrt{15}$ square inches. The length of the third median, in inches, is:


 $\textbf{(A)}\ 4 \qquad  \textbf{(B)}\ 3\sqrt{3} \qquad  \textbf{(C)}\ 3\sqrt{6} \qquad  \textbf{(D)}\ 6\sqrt{3} \qquad  \textbf{(E)}\ 6\sqrt{6}$

 \textbf{Answer: }C



\textbf{Solution}

By the area formula:
 \[A = \frac43\sqrt{s(s-m_1)(s-m_2)(s-m_3)}\]Where $s = \frac{m_1+m_2+m_3}{2}$.Plugging in the numbers:
 \[3\sqrt{15} = \frac43\sqrt{\frac{9+m}2\cdot\frac{9-m}2\cdot\frac{m+3}2\cdot\frac{m-3}2}\]Simplifying and squaring both sides:
 \[1215 = (9-m)(9+m)(m+3)(m-3)\]Now, we can just plug in the answer choices and find that $\boxed{3\sqrt6}$ works.

 



\textbf{Solution 2}

We connect all the medians of the triangle. We know that the ratio of the vertice to orthocenter / median is in a $2:3$ ratio. For convenience, call the orthocenter $O$, and label the triangle $\triangle ABC$ such that the median from $A$ to $BC$ is 3 and the median from $B$ to $AC$ is 6. Then, $BO$ is 4 and $AO$ is 2. Let us call the portion of the third median that goes from $O$ to $AB$ have length $x$, and and $OC$ have length $2x$. Note that the median from $O$ to $AB$ in $\triangle AOB$ is equal to $x$. 

 Note that the medians split the triangle into 6 triangles of equal area, so $\triangle AOB$ has area equal to $\frac{1}{3}$ of $\triangle ABC = \sqrt{15}$. Let $AB=2s$. Using Herons*, we get:

 
 \[15=(3+s)(s-1)(s+1)(3-s)\]
 \[=(9-s^2)(s^2-1)\]

 We can see that $s^2=4$, meaning that $s$ is 2 and $AB=4$. We can then use Stewart's to find the length of the median from $O$, since we know the median cuts $AB$ into segments each of length $2$. We get:

 
 \[(2^2)(4)+4x^2=(4^2)(2)+(2^2)(2)\]
 \[16+4x^2=32+8\]
 \[4x^2=24\]
 \[x^2=6\]
 \[x=\sqrt{6}\]

 Since the length of the actual median from $C$ is equal to $3$, we have that the answer is $\boxed{3\sqrt{6}}$.

 
            \begin{itemize}
        \item If anyone has a better method of either finding $AB$ or the median of $AOB$ from $O$, please feel free to edit
            \end{itemize}
            


	\item (\textit{AHSME 1962}) The limiting sum of the infinite series

 \[\frac{1}{10} + \frac{2}{10^2} + \frac{3}{10^3} + \cdots\]
whose $n$-th term is $n/10^n$ is:


 $\textbf{(A)}\ \frac{1}{9} \qquad  \textbf{(B)}\ \frac{10}{81} \qquad  \textbf{(C)}\ \frac{1}{8} \qquad  \textbf{(D)}\ \frac{17}{72} \qquad  \textbf{(E)}\ \text{larger than any finite quantity}$

 \textbf{Answer: }B



\textbf{Solution}

The series can be written as the following:

 $\frac{1}{10} + \frac{1}{10^2} + \frac{1}{10^3} + ...$
$+ \frac{1}{10^2} + \frac{1}{10^3} + \frac{1}{10^4} + ...$
$+ \frac{1}{10^3} + \frac{1}{10^4} + \frac{1}{10^5} + ...$

 and so on.

 by using the formula for infinite geometric series $(\frac{a}{1-r})$,

 We can get $\frac{\frac{1}{10}}{1-\frac{1}{10}}$ $+$ $\frac{\frac{1}{10^2}}{1-\frac{1}{10}}$ $+$ $\frac{\frac{1}{10^3}}{1-\frac{1}{10}}$ $+$ ...Since they all have common denominators, we get $\frac{(\frac{1}{10} + \frac{1}{10^2} + \frac{1}{10^3})}{\frac{9}{10}}$.Using the infinite series formula again, we get $\frac{\frac{\frac{1}{10}}{1-\frac{1}{10}}}{\frac{9}{10}}$ $=$ $\frac{\frac{\frac{1}{10}}{\frac{9}{10}}}{\frac{9}{10}}$ $=$ $\frac{\frac{1}{9}}{\frac{9}{10}}$ $=$ $\boxed{ (B) \frac{10}{81}}$




\textbf{Solution 2}

So.. we have the sum to be $\frac{1}{10}+\frac{2}{100}+\frac{3}{1000}$...Notice that this can be written as $\frac{1}{10}+\frac{0.2}{10}+\frac{0.03}{10}+\frac{0.004}{10}$.Now, it is trivial that the new fraction we seek is $\frac{1.234567891011......}{10}$

 Testing the answer choices, we see that $\boxed{B}$ is the correct answer.

 



\textbf{Solution 3}

Let
 \[S = \frac{1}{10} + \frac{2}{10^2} + \frac{3}{10^3} + ...\]Then
 \[10S = 1 + \frac{2}{10} + \frac{3}{10^2} + \frac{4}{10^3} + ...\]Subtracting $1S$ from $10S$, we got:
 \begin{align*} 9S &= 1 + \frac{1}{10} + \frac{1}{10^2} + \frac{1}{10^3} + ... \\ &= \frac{1}{1-\frac{1}{10}} = \frac{10}{9} \\ S &= \frac{10}{81} \end{align*}
Therefore, the answer is $\boxed{(B) \frac{10}{81}}$. -nullptr07

 


\end{enumerate}